%!TEX root = forallxyyc.tex

\chapter{Notação simbólica}
\label{app.notation}

\section{Nomenclatura alternativa}

\paragraph{Lógica verofuncional} A LVF recebe outros nomes. Às vezes, é chamada de \emph{lógica sentencial}, porque lida fundamentalmente com sentenças. Às vezes, é chamada de \emph{lógica proposicional}, pela ideia de que lida fundamentalmente com proposições. Optamos por \emph{lógica verofuncional} para enfatizar que ela trata apenas de atribuições de verdade e falsidade a sentenças e que todos os seus conectivos são verofuncionais.

\paragraph{Lógica de primeira ordem} A LPO recebe outros nomes. Às vezes, é chamada de \emph{lógica de predicados}, porque permite aplicar predicados a objetos. Às vezes, é chamada de \emph{lógica quantificacional}, porque faz uso de quantificadores.

\paragraph{Fórmulas} Alguns textos chamam as fórmulas de \emph{fórmulas bem-formadas}. Como \textquotedblleft fórmula bem-formada\textquotedblright{} é uma expressão longa e incômoda, esses textos a abreviam como \emph{wff}. Isso é ao mesmo tempo bárbaro e desnecessário (tais textos não admitem \textquotedblleft fórmulas malformadas\textquotedblright{}). Optamos por \textquotedblleft fórmula\textquotedblright{}.

Em \cref{s:TFLSentences}, definimos as \emph{sentenças} da LVF. Às vezes, elas também são chamadas de \textquotedblleft fórmulas\textquotedblright{} (ou \textquotedblleft fórmulas bem-formadas\textquotedblright{}), pois, na LVF, ao contrário da LPO, não há distinção entre uma fórmula e uma sentença.

\paragraph{Valorações} Alguns textos chamam as valorações de \emph{atribuições de verdade} ou \emph{atribuições de valores de verdade}.

\paragraph{Predicados de $n$ lugares} Escolhemos chamar os predicados de \textquotedblleft de um lugar\textquotedblright{}, \textquotedblleft de dois lugares\textquotedblright{}, \textquotedblleft de três lugares\textquotedblright{} etc. Outros textos os chamam, respectivamente, de \textquotedblleft monádicos\textquotedblright{}, \textquotedblleft diádicos\textquotedblright{}, \textquotedblleft triádicos\textquotedblright{} etc. Outros ainda os chamam de \textquotedblleft unários\textquotedblright{}, \textquotedblleft binários\textquotedblright{}, \textquotedblleft ternários\textquotedblright{} etc.

\paragraph{Nomes} Na LPO, usamos \textquotedblleft $a$\textquotedblright{}, \textquotedblleft $b$\textquotedblright{} e \textquotedblleft $c$\textquotedblright{} como nomes. Alguns textos os chamam de \textquotedblleft constantes\textquotedblright{}. Outros textos não marcam nenhuma diferença entre nomes e variáveis na sintaxe. Esses textos se concentram simplesmente em saber se o símbolo ocorre \emph{ligado} ou \emph{livre}.

\paragraph{Domínios} Alguns textos descrevem um domínio como um \textquotedblleft domínio do discurso\textquotedblright{} ou um \textquotedblleft universo do discurso\textquotedblright{}.

\section{Símbolos alternativos}
Na história da lógica formal, diferentes símbolos foram usados em diferentes épocas e por diferentes autores. Muitas vezes, os autores eram obrigados a usar a notação que suas gráficas conseguiam compor. Este apêndice apresenta alguns símbolos comuns, para que você possa reconhecê-los caso os encontre em um artigo ou em outro livro.

\paragraph{Negação} Dois símbolos usados com frequência são o \emph{gancho}, \textquotedblleft $\neg$\textquotedblright{}, e o \emph{traço ondulado}, ou \emph{til}, \textquotedblleft $\sim$\textquotedblright{}. Em alguns sistemas formais mais avançados, é necessário distinguir entre dois tipos de negação; essa distinção às vezes é representada usando tanto \textquotedblleft $\neg$\textquotedblright{} quanto~\textquotedblleft $\sim$\textquotedblright{}. Textos mais antigos às vezes indicam a negação por uma linha sobre a fórmula negada, por exemplo, \textquotedblleft $\overline{A \eand B}$\textquotedblright{}, ou por um sinal de menos~\textquotedblleft $-$\textquotedblright{}. Muitos textos usam \textquotedblleft $x \neq y$\textquotedblright{} como abreviação de \textquotedblleft $\enot x = y$\textquotedblright{}.

\paragraph{Disjunção} O símbolo \textquotedblleft $\vee$\textquotedblright{} é normalmente usado para simbolizar a disjunção inclusiva. Em alguns textos da chamada tradição algébrica, a disjunção é escrita como uma adição~\textquotedblleft $+$\textquotedblright{}. Em muitas linguagens de programação, usa-se uma barra vertical~\textquotedblleft \verb+|+\textquotedblright{} (ou duas: \textquotedblleft \verb+||+\textquotedblright{}). (Isso pode ser confundido com a barra de Sheffer de \cref{c:FunctionalCompleteness}, que é normalmente escrita como~\textquotedblleft $\mid$\textquotedblright{}.)

\paragraph{Conjunção}
A conjunção é frequentemente simbolizada pelo \emph{e comercial}, \textquotedblleft {\&}\textquotedblright{}. O e comercial é uma forma decorativa da palavra latina \textquotedblleft et\textquotedblright{}, que significa \textquotedblleft e\textquotedblright{}. (Sua etimologia ainda aparece em certas fontes, especialmente nas fontes em itálico; assim, um e comercial em itálico pode aparecer como \textquotedblleft%
\iflatexml
\<span class="ltx\_text" style="font-family:Palatino; font-style:italic;">\&\</span>
\else
\emph{\&}
\fi
\textquotedblright{}.) Esse símbolo é usado comumente na escrita natural em inglês (por exemplo, \textquotedblleft Smith \& Sons\textquotedblright{}) e, embora seja uma escolha natural, muitos lógicos usam um símbolo diferente para evitar confusão entre a linguagem-objeto e a metalinguagem: como símbolo de um sistema formal, o e comercial não é a palavra inglesa \textquotedblleft \&\textquotedblright{}. A escolha mais comum atualmente é \textquotedblleft $\wedge$\textquotedblright{}, que corresponde ao símbolo usado para a disjunção. Às vezes, usa-se um único ponto, \textquotedblleft {\scriptsize\textbullet}\textquotedblright{}. Em alguns textos mais antigos, não há símbolo para a conjunção; \textquotedblleft $A$ e $B$\textquotedblright{} é simplesmente escrito como \textquotedblleft $AB$\textquotedblright{}.

\paragraph{Condicional material} Há dois símbolos comuns para o condicional material: a \emph{seta}, \textquotedblleft $\rightarrow$\textquotedblright{}, e a \emph{ferradura}, \textquotedblleft $\supset$\textquotedblright{}.

\paragraph{Bicondicional material} A \emph{seta de duas pontas}, \textquotedblleft $\leftrightarrow$\textquotedblright{}, é usada nos sistemas que usam a seta para representar o condicional material. Os sistemas que usam a ferradura para o condicional normalmente usam a \emph{barra tripla}, \textquotedblleft $\equiv$\textquotedblright{}, para o bicondicional.

\paragraph{Quantificadores} O quantificador universal é normalmente simbolizado como um \textquotedblleft A\textquotedblright{} girado, e o quantificador existencial como um \textquotedblleft E\textquotedblright{} girado. Em alguns textos, não há um símbolo separado para o quantificador universal. Em vez disso, a variável é simplesmente escrita entre parênteses diante da fórmula que ela liga. Por exemplo, eles podem escrever \textquotedblleft $(x)\atom{P}{x}$\textquotedblright{} onde escreveríamos \textquotedblleft $\forall x\, \atom{P}{x}$\textquotedblright{}. Alguns textos também usam versões grandes dos nossos conectivos de conjunção e disjunção, isto é, \textquotedblleft $\bigwedge$\textquotedblright{} e \textquotedblleft $\bigvee$\textquotedblright{}, para os quantificadores universal e existencial, respectivamente. (Às vezes, a variável é colocada como subscrito desses símbolos, ou até mesmo abaixo deles.)

\bigskip

As tipografias alternativas são resumidas abaixo:

\begin{center}
\begin{tabular}{rl}
negação & $\neg$, $\sim$, $-$, $\overline{A}$\\
conjunção & $\wedge$, $\&$, {\scriptsize\textbullet}\\
disjunção & $\vee$, $+$, \verb+|+, \verb+||+\\
condicional & $\rightarrow$, $\supset$\\
bicondicional & $\leftrightarrow$, $\equiv$\\
quantificador universal & $\forall x$, $(x)$, $\bigwedge_x$\\
quantificador existencial &$\exists x$, $\bigvee_x$
\end{tabular}
\end{center}





\section{Notação polonesa}

Esta seção apresenta brevemente a lógica sentencial em notação polonesa, um sistema de notação introduzido no final da década de 1920 pelo lógico polonês Jan {\L}ukasiewicz.

Letras minúsculas são usadas como letras sentenciais. A letra maiúscula $N$ é usada para a negação. $A$ é usada para a disjunção, $K$ para a conjunção, $C$ para o condicional e $E$ para o bicondicional. (\textquotedblleft $A$\textquotedblright{} é uma abreviação de alternância, outro nome para a disjunção lógica. \textquotedblleft $E$\textquotedblright{} é uma abreviação de equivalência.)

Na notação polonesa, um conectivo binário é escrito \emph{antes} das duas sentenças que ele conecta. Por exemplo, a sentença $A\eand B$ da LVF seria escrita como $Kab$ na notação polonesa.

As sentenças $\enot A\eif B$ e $\enot (A\eif B)$ são muito diferentes; o operador lógico principal da primeira é o condicional, mas o conectivo principal da segunda é a negação. Na LVF, mostramos isso colocando parênteses ao redor do condicional na segunda sentença. Na notação polonesa, parênteses nunca são necessários. O conectivo mais à esquerda é sempre o conectivo principal. A primeira sentença seria simplesmente escrita como $CNab$ e a segunda como $NCab$.

Essa característica da notação polonesa significa que é possível avaliar sentenças simplesmente percorrendo os símbolos da direita para a esquerda. Se você estivesse construindo uma tabela-verdade para $NKab$, por exemplo, primeiro consideraria os valores de verdade atribuídos a $b$ e $a$, depois consideraria a conjunção entre eles e, por fim, negaria o resultado. A regra geral para decidir o que avaliar em seguida na LVF não é nem de longe tão simples. Na LVF, a tabela-verdade de $\enot(A\eand B)$ exige olhar para $A$ e $B$, depois olhar para a conjunção no meio da sentença e, em seguida, para a negação no início da sentença. Como a ordem das operações pode ser especificada de maneira mais mecânica na notação polonesa, variantes da notação polonesa são usadas como estrutura interna de muitas linguagens de programação de computadores.
